\documentclass[11pt]{article}
\usepackage[T1]{fontenc}
\usepackage[margin=0.85in]{geometry}
\usepackage{amsmath,amssymb}
\usepackage{booktabs,array,graphicx}
\usepackage[expansion=false]{microtype}
\setlength{\emergencystretch}{3em}
\setlength{\parindent}{0pt}
\setlength{\parskip}{0.45em}

\title{Supplementary material for\\
Two-stage MILP scheduling and structure-aware learning-enhanced branch-and-cut for dual-source mixed-flow rotating-chamber cluster tools}
\author{Anonymous author(s)}
\date{}

\begin{document}
\maketitle

\section{Computational evidence design}

The computational study separates formulation validation, nominal method
comparison, mechanism-specific controls, manufacturing-factor analysis,
parameter sensitivity, and distribution-shift evaluation. Table~\ref{tab:suites}
gives the complete prespecified design. A run is identified by the instance,
method, solver seed, training seed where applicable, and experimental profile.
All observations are retained, including time-limited runs.

\begin{table}[htbp]
\centering
\caption{Prespecified computational suites.}
\label{tab:suites}
\begin{tabular}{p{0.18\linewidth}p{0.52\linewidth}r}
\toprule
Suite & Design & Runs \\
\midrule
Validation & Four small physical cases, one exact solve each & 4 \\
Nominal & 20 instances; two solver baselines with five solver seeds and five learned configurations with five solver and five training seeds & 2700 \\
Two-stage stability & 20 instances, three refinement profiles, five solver seeds, and five training seeds & 1500 \\
Factorial DOE & $4\times5\times3$ wafer-count, recipe-composition, and processing-scale cells with five solver seeds & 300 \\
SPBS & 48 instances, automatic versus no warm start, ten matched solver seeds & 960 \\
Sensitivity & Eight one-factor settings with ten solver seeds & 80 \\
Distribution shift & Nine out-of-distribution instances; two solver baselines and five learned configurations with matched solver and training realizations & 729 \\
\midrule
Total & Excluding the four validation cases & 6269 \\
\bottomrule
\end{tabular}
\end{table}

All 6269 formal benchmark runs and four validation runs completed without
execution failure. Every saved incumbent used in the reported statistics was
checked independently against the corresponding mathematical instance. The
distribution-shift suite contains 729 retained time-limited runs; every run
returns an independently valid incumbent.

\section{Indexed MILP specification}

This section records the principal indexed constraint families used to
instantiate the finite-horizon model; it is a compact mathematical map of the
implemented formulation rather than a claim that every generated row is
printed here. Let $\mathcal W$, $\mathcal P$, $\mathcal C$, $\mathcal K$,
$\mathcal O$, and $\mathcal R$ denote product wafers, same-mode product-pair
groups, rotary chambers, candidate recipe slots, operations, and unit-capacity
resources. A slot $k\in\mathcal K(p)$ is a chamber--batch--side tuple in
$4\times1$ mode or a chamber--chain-position tuple in $2\times2$ mode.
$\mathcal A$ is the set of technological predecessor arcs, and
$\mathcal O_r$ contains operations that may use resource $r$. For each chamber
$c$, $\mathcal V_c=(v_{c1},\ldots,v_{c n_c})$ is the ordered set of candidate
processing cycles and $\mathcal E_c=\{1,\ldots,E_c\}$ is the ordered set of
potential cleaning epochs. The parameter $H$ is the maximum number of active
processing cycles permitted between two cleans. Furthermore,
$\mathcal J_c^{\rm DW}$ contains dummy-wafer use demands and $\mathcal T_c$
contains the reusable physical vacuum-zone dummy wafers.

The binary variables are assignment $x_{pk}$, position activation $u_k$,
operation activation $b_o$, cycle-to-epoch assignment $a_{vce}$, epoch use
$z_{ce}$, dummy identity $g_{j\tau}$, and resource order $y^r_{oo'}$. Here
$b_v$ is the activation variable $b_o$ for the processing operation associated
with cycle $v\in\mathcal V_c$, and $z_{ce}=1$ means that epoch $e$ of chamber
$c$ contains at least one active cycle. Likewise, $b_j$ below abbreviates the
activation of dummy-use demand $j$. Continuous variables $s_o,e_o$ are event
start and completion times.

Pairs belonging to one recipe mode have identical processing data and their
identifiers carry no physical meaning. They are therefore ordered by identifier
and mapped to the compatible slots in a fixed interleaved chamber order. Define
the canonical incidence
\begin{equation}
\bar{x}_{pk}=\mathbf 1\{k=k^{\rm can}(p)\},\qquad
x_{pk}=\bar{x}_{pk}\quad(p\in\mathcal P,\ k\in\mathcal K(p)).
\label{eq:canonical-incidence}
\end{equation}
Equation~\ref{eq:canonical-incidence} is a same-mode symmetry reduction: it
chooses one representative from
each orbit of pair-label permutations and does not remove a physically distinct
schedule under the modeled identical-pair data. The odd product/dummy tail, if
present, occupies the last compatible $2\times2$ position. The resulting
fixed-duration, activation, and assignment relations are
\begin{align}
&\sum_{k\in\mathcal K(p)}x_{pk}=1 &&(p\in\mathcal P),\\
&x_{pk}\le u_k,\quad u_{k+1}\le u_k,\quad
b_o=\chi_o(x,u,a) &&(o\in\mathcal O),\\
&b_{o'}\le b_o &&((o,o')\in\mathcal A),\\
&0\le s_o\le Mb_o,\quad e_o=s_o+d_ob_o &&(o\in\mathcal O),
\end{align}
where the binary affine incidence map $\chi_o$ activates exactly the events
belonging to the selected $4\times1$ side, $2\times2$ chain position, or
cleaning cycle. The predecessor inequality makes operation activation
consistent along every generated technological chain: an active successor
cannot have an inactive required predecessor. The route logic is generated
from the following active event chains:
\begin{center}
\begin{tabular}{p{0.16\linewidth}p{0.76\linewidth}}
\toprule
Structure & Ordered active events \\
\midrule
Product route & load port $\rightarrow$ ATR $\rightarrow$ aligner
$\rightarrow$ upper load lock $\rightarrow$ VTR load $\rightarrow$ chamber
process $\rightarrow$ VTR unload $\rightarrow$ lower load lock
$\rightarrow$ ATR $\rightarrow$ load port \\
$4\times1$ & front load $\rightarrow$ rotate/back load $\rightarrow$
two-side process $\rightarrow$ back unload $\rightarrow$ rotate/front unload \\
$2\times2$ & dummy/product head $\rightarrow$ adjacent product bridges
$\rightarrow$ product/dummy tail; consecutive positions share the prescribed
product pair and may split only at a cleaning boundary \\
Cleaning & front dummy load $\rightarrow$ back dummy load $\rightarrow$
clean process $\rightarrow$ back unload $\rightarrow$ front unload \\
\bottomrule
\end{tabular}
\end{center}

For every active technological arc and every unordered pair competing for a
resource, timing feasibility is imposed by
\begin{align}
s_{o'}&\ge e_o+\delta_{oo'}-M(2-b_o-b_{o'})
&&(o,o')\in\mathcal A,\\
s_{o'}&\ge e_o+\rho^r_{oo'}-M(1-y^r_{oo'})
-M(2-b_o-b_{o'}),\\
s_o&\ge e_{o'}+\rho^r_{o'o}-My^r_{oo'}
-M(2-b_o-b_{o'})
&&\{o,o'\}\subseteq\mathcal O_r.
\end{align}
The setup $\rho^r_{oo'}$ includes the empty reposition time from the endpoint
of $o$ to the origin of $o'$. These disjunctions cover the atmospheric and
vacuum robots, aligner, individual load-lock slots, rotary chambers, and
physical dummy-wafer identities.

Dummy identity and cleaning are linked by
\begin{align}
&\sum_{\tau\in\mathcal T_c}g_{j\tau}=b_j
&&(j\in\mathcal J_c^{\rm DW}),\\
&\sum_{e\in\mathcal E_c}a_{vce}=b_v,\quad
a_{vce}\le z_{ce},\quad
z_{ce}\le\sum_{v\in\mathcal V_c}a_{vce},\\
&\sum_{v\in\mathcal V_c}a_{vce}\le H z_{ce},\quad
z_{c,e+1}\le z_{ce},\\
&b_{v_{c,q+1}}\le b_{v_{cq}},\\
&\sum_{e\in\mathcal E_c}e\,a_{v_{cq},c,e}
\le\sum_{e\in\mathcal E_c}e\,a_{v_{c,q+1},c,e}
+E_c(1-b_{v_{c,q+1}})
&& (q=1,\ldots,n_c-1).
\end{align}
If $g_{j\tau}=g_{j'\tau}=1$, the same resource disjunction is activated for
dummy uses $j$ and $j'$, preventing overlapping reuse. Epoch transitions
activate the intervening dummy-only cleaning route, with its activation equal
to $z_{c,e+1}$. The two last inequalities impose prefix activation and
nondecreasing epoch indices on consecutive active cycles; the term
$E_c(1-b_{v_{c,q+1}})$ prevents an inactive suffix from falsely constraining
its active predecessor. The adjacency equations in $\chi_o$ permit a
$2\times2$ split only at an activated epoch transition.

Residence and completion constraints are
\begin{align}
0&\le s_{\operatorname{out}(w,r)}
-e_{\operatorname{in}(w,r)}\le\overline R_{wr},\\
C_{\max}&\ge e_w^{\rm return} &&(w\in\mathcal W).
\end{align}
Stage one minimizes $C_{\max}$. Stage two fixes the retained routing and
resource-order integers from its stage-one incumbent (double-gripper choices
remain free), adds $C_{\max}\le\widehat C+10^{-6}$, and minimizes
\begin{equation}
\Phi=L_{\rm sum}+10L_{\max}+25D_{\rm cadence}+0E_{\rm cap}.
\end{equation}
Let $w_i\ge0$ denote a model-native wait or capacity-adjusted idle term,
$i\in\mathcal L$, let $\bar{w}_i$ be its optional soft target, and let
$h_{mce\ell}$ be the $\ell$th active processing start of mode $m$, chamber
$c$, and epoch $e$. The max, absolute-value, and optional soft-cap terms are
linearized by
\begin{align}
L_{\rm sum}&=\sum_{i\in\mathcal L}w_i,&
L_{\max}&\ge w_i &&(i\in\mathcal L),\\
\Delta_{mce\ell}&=(h_{mce,\ell+1}-h_{mce,\ell})
 -(h_{mce,\ell}-h_{mce,\ell-1}),\\
d_{mce\ell}&\ge \Delta_{mce\ell},&
d_{mce\ell}&\ge-\Delta_{mce\ell},&
D_{\rm cadence}&=\sum_{m,c,e,\ell\ge2}d_{mce\ell},\\
\xi_i&\ge w_i-\bar{w}_i,&
\xi_i&\ge0,&
E_{\rm cap}&=\sum_{i\in\mathcal L}\xi_i.
\end{align}
The symbol $\xi_i$ is reserved for cap excess and is distinct from the epoch
indicator $z_{ce}$. A positive coefficient makes each epigraph variable tight
at optimum. The formal profile disables the soft target, so its effective cap
weight is zero and only the first three terms enter the reported objective.
Hence stage two preserves the stage-one
makespan within numerical tolerance; if stage one is time-limited, this
statement applies to its retained incumbent rather than to an unproved global
optimum.

For the deterministic offline problem, a scheduling deadlock is a circular
wait among active material-transfer or processing operations in which every
operation awaits a resource or predecessor held ahead of it. Such a wait would
be a directed cycle in the combined technological and selected resource-order
graph. Summing the positive separation inequalities around that cycle gives
$s_o\ge s_o+\epsilon$ for some $\epsilon>0$, a contradiction; terminal
constraints also require every product and reusable dummy route to finish.
Thus every feasible solution is a deadlock-free complete offline trajectory
for the modeled resources, durations, and availability assumptions. This is
not a runtime deadlock-recovery guarantee under equipment faults, stochastic
delays, operator intervention, or any resource state omitted from the model.

\section{Independent validation and numerical tolerance}

The four small validation cases exercise pure $4\times1$, pure $2\times2$,
and two mixed configurations. Their saved schedules are reloaded, the reported
primary decisions are fixed, and the active physical model is checked
independently of the learning policy. The same solution-level procedure is
then applied to every saved benchmark incumbent. It covers route and stage
order, chamber and slot exclusivity, robot and load-lock use, recipe grouping,
cleaning, residence limits, and vacuum-zone dummy-wafer capacity. Continuous
residuals use the solver tolerance $10^{-6}$ plus a $10^{-12}$ serialization
allowance, i.e., an effective threshold of
$1.000001\times10^{-6}$. Binary and integer requirements are not relaxed.
This allowance only accounts for the decimal text round trip of an otherwise
accepted solver solution.

The validation claim is deliberately bounded: it establishes consistency
between the generated instances, serialized solutions, and implemented
constraints for the tested cases. It is not field validation of unmodeled
equipment behavior, a proof of correctness for every possible parameterization,
or evidence of online recovery from faults and stochastic disturbances.

\section{Event-time reconstruction of schedule stability}

Schedule-stability metrics are reconstructed from physical event times rather
than auxiliary objective variables. This makes the definition identical under
all three refinement profiles. For product wafer $w$, define six nonnegative
route waits:
\begin{align}
q_{w,1}&=[s^{\mathrm{PM}}_w-e^{\mathrm{VTRload}}_w]_+, &
q_{w,2}&=[s^{\mathrm{VTRunload}}_w-e^{\mathrm{PM}}_w]_+,\\
q_{w,3}&=[s^{\mathrm{AL}}_w-e^{\mathrm{ATR,LP\to AL}}_w]_+, &
q_{w,4}&=[s^{\mathrm{ATR,AL\to LL_u}}_w-e^{\mathrm{AL}}_w]_+,\\
q_{w,5}&=[s^{\mathrm{VTRload}}_w-e^{\mathrm{LL_u}}_w]_+, &
q_{w,6}&=[s^{\mathrm{ATR,LL_l\to LP}}_w-e^{\mathrm{LL_l}}_w]_+,
\end{align}
where $[x]_+=\max(0,x)$ and $s$ and $e$ denote stage start and end times.
The route-total and route-maximum metrics are
\begin{equation}
Q_{\mathrm{total}}=\sum_w\sum_{k=1}^{6}q_{w,k},\qquad
Q_{\max}=\max_{w,k}q_{w,k}.
\end{equation}

Processing cadence is reconstructed from the ordered start times of active
full-flow and mixed-flow chamber batches. For adjacent-start intervals
$d_1,\ldots,d_m$, cadence variation is
$\operatorname{CV}(d)=\operatorname{sd}(d)/\operatorname{mean}(d)$.
The cadence-deviation sum is the sum of the absolute deviations used by the
linear refinement objective after the discrete batch structure has been fixed.
Table~\ref{tab:stability-profiles} reports the resulting instance-first means.

\begin{table}[htbp]
\centering
\caption{Instance-first means in the matched, equal-total-budget
stability-profile experiment.}
\label{tab:stability-profiles}
\begin{tabular}{lrrrr}
\toprule
Profile & Route-total & Route-maximum & Cadence CV & Cadence deviation \\
\midrule
Full & 2596.01 & 296.33 & \textbf{0.130} & \textbf{22.04} \\
Quadratic objective (linear off) & \textbf{2469.47} & \textbf{293.27} & 0.134 & 60.18 \\
Stage2-off & 5398.60 & 952.57 & 0.216 & 370.41 \\
\bottomrule
\end{tabular}
\end{table}

Relative to stage2-off, the full profile reduces route-total wait by 51.91\%,
route-maximum wait by 68.89\%, cadence CV by 39.70\%, and cadence deviation by
94.05\%; all four comparisons remain significant after Holm adjustment. The
linear objective also reduces cadence deviation by 63.38\% relative to the
quadratic profile. Its route-total, route-maximum, and cadence-CV differences
from the quadratic profile are not Holm-adjusted significant. These estimates
compare complete stability profiles under the same 600-s total budget, not a
pure stage-two intervention applied to an identical stage-one incumbent:
stage2-off can spend the full budget on the primary solve, whereas the two
refinement profiles reserve up to 60 s. The results therefore identify the
operational effect of the prespecified full two-stage configuration under a
common total budget; they do not isolate a stage-two causal effect conditional
on an equal primary search horizon.

\section{Statistical protocol}

The manufacturing instance is the unit of inference. Repeated solver and
training realizations are first averaged within each instance and method or
profile. Mean differences are reported with a nonparametric 95\% bootstrap
interval. Paired two-sided Wilcoxon signed-rank tests use the same instance
pairs, and all planned contrasts within an analysis family are adjusted by
the Holm procedure. Failure, incumbent, and optimality indicators are retained
as binary outcomes. Primal--dual integral (PDI) is the primary anytime measure;
lower values are better. Time-limited runs remain in PDI, time, incumbent, and
failure summaries.

Specifically, if $G(t)$ is SCIP's normalized primal--dual gap at time $t$,
\begin{equation}
\operatorname{PDI}(T)=\int_0^T G(t)\,dt .
\end{equation}
Terminal gap is $G(T)$. Reported solution time is the sum of SCIP-reported
primary and conditional-refinement solving times. It excludes model loading,
solution serialization, and offline SPBS construction. The paired SPBS suite
therefore estimates the downstream solve effect of loading a precomputed
start; it is not an end-to-end comparison that charges start construction.

\section{Solver, learning, and decoding configuration}

All nominal methods receive the same precomputed, independently verified SPBS
incumbent. Consequently, the 2700-run comparison isolates cut-selection and
decoding policies under common primal guidance; it does not estimate an
SPBS--policy interaction. The seven configurations are:
\begin{description}
\item[SCIP default] Native SCIP cut selection with the shared SPBS incumbent.
\item[ACS] A four-score SCIP selector with direction, efficacy, integer-support,
and objective-parallelism weights $(0.25,0,0.50,0.25)$.
\item[HEM-13D] The adapted 13-feature hierarchical policy with greedy decoding.
\item[HEM-13D+Beam] The identical trained HEM policy parameters with width-three beam decoding.
\item[23D Features Only] The reconstructed 23-feature policy with greedy
decoding and no role-based reranking.
\item[Structure+Greedy] The trained proposed policy parameters and role-based completion with
greedy decoding.
\item[RDMCT-HBS] The same trained proposed policy parameters and completion rule with
width-three beam decoding.
\end{description}

Each candidate cut has 13 generic attributes: objective parallelism, efficacy,
support, integer support, normalized violation, and the mean, maximum, minimum,
and standard deviation of cut and objective coefficients. The ten added
descriptors use only solver-visible variable types and optimization roles:
coefficient-mass shares on SCIP binary, general-integer (including implied
integer), nonzero-objective continuous, zero-objective auxiliary-continuous,
and other variables, followed by bounded-variable mass, positive-coefficient
mass, discrete--continuous coupling, dominant-family mass, and normalized
family entropy. They are computed from SCIP type, objective, bound, column, and
row-coefficient metadata; formulation-specific variable names are never
parsed. Candidates are screened deterministically by descending efficacy,
with original index as the tie breaker, to at most 128 rows; at most 16 rows
are selected.

The 23-dimensional input passes through layer normalization, separate
$13$- and $10$-dimensional projections, and a learned residual gate before a
one-layer LSTM encoder and pointer-attention decoder. The recurrent hidden
dimension is 128, with one attention glimpse and tanh exploration coefficient
5. A high-level policy chooses the selection fraction; a low-level pointer
policy with an explicit end token generates an ordered, nonrepeating sequence.
Beam width is three, ranked by cumulative joint log probability without length
normalization. Structural completion retains at most
$\lceil0.25k\rceil$ policy anchors for target cardinality $k$ and greedily
fills from a pool of at most $2k$ candidates. Its quality, policy, role
coverage, and representativeness weights are $0.35$, $0.20$, $0.25$, and
$0.20$. The associated $(1-1/e)$ property applies only to the residual
monotone-submodular completion problem with fixed anchors and pool, not to
global cut selection.

Training uses five training LPs and five independent training seeds. For each
seed, two workers jointly collect eight complete trajectories per epoch for 60
epochs with stochastic decoding; each trajectory has a 30-s/1024-MB SCIP
budget, including at most 5 s of conditional refinement and one separation
round at the root. A central synchronous Adam update uses an
exponential-moving-average return baseline with coefficient $0.9$. Low- and
high-level policy learning rates are $10^{-4}$ and $5\times10^{-4}$; the
gradient norm is capped at 2, and both learning rates are multiplied by 0.96
every ten epochs. The return is negative SCIP PDI, with neither feature nor
reward normalization. This is an A3C-inspired parallel hierarchical REINFORCE
backbone with no learned critic and no asynchronous parameter update.

The terminal training epoch is prespecified rather than selected on benchmark
performance. For each method family and seed, the retained epoch-60 parameter
set must match the training split and seed and must have the required feature
dimension and PDI reward. There is one admissible parameter set for each
family--seed pair; no benchmark or test outcome enters parameter selection.

\section{SPBS construction and acceptance}

SPBS first fixes the canonical assignment, activation, and cleaning-epoch
core in a temporary zero-objective feasibility MIP. The profiles then add or
release load-lock-slot and resource-order decisions in the deterministic
sequence shown in Table~\ref{tab:spbs-profiles}. Fractions are shares of the
configured SPBS construction limit; every nonfinal attempt receives the
smaller of its share and the remaining time, with a 5-s minimum when time
remains, and the final profile receives the remainder. Each attempt stops at
its first feasible solution.

\begin{table}[htbp]
\centering
\caption{Deterministic SPBS release hierarchy for instances below 32 product
wafers. The core comprises canonical assignment, activation, and cleaning
epochs.}
\label{tab:spbs-profiles}
\begin{tabular}{p{0.20\linewidth}p{0.58\linewidth}r}
\toprule
Profile & Fixed decisions & Limit share \\
\midrule
Full & Core, load-lock slots, ATR/aligner, load-lock, and VTR orders & 0.10 \\
Relaxed-VTR & Core, slots, ATR/aligner and load-lock orders & 0.15 \\
Relaxed-LL & Core, slots and ATR/aligner orders & 0.45 \\
Fixed-slots & Core and load-lock slots & 0.15 \\
Structure-only & Core only & remainder \\
\bottomrule
\end{tabular}
\end{table}

For at least 32 product wafers without a positive wait cap, the sequence is
fixed-slots (0.50), structure-only (0.25), and an unrestricted final fallback.
With a positive wait cap, it is structure-only (0.60) followed by the
unrestricted fallback. The unrestricted profile removes even the structural
fixings and uses only the remaining feasibility-search budget.

The selected physical schedule is reconstructed with secondary variables free
in the full current model and is accepted only after SCIP's complete original-
space feasibility check. The production MILP receives only the accepted
incumbent; no temporary fixing survives. Each cached start records the accepted
profile and construction attempts and is reusable only when its full-model
scope and exact LP SHA-256 match the current instance.

\section{Hardware, software, and budgets}

Experiments ran under Windows 11 Pro on an Intel Core Ultra 7 265KF
(20 cores), 31.7 GB RAM, and an NVIDIA GeForce RTX 5060 Ti with 8 GB memory.
The evaluation environment uses Python 3.12.12, SCIP 10.0.1 through
PySCIPOpt 6.1.0, and PyTorch 2.11.0+cu130. The code does not impose a solver thread count, so no
single-thread assumption is used in the analysis.
Table~\ref{tab:budgets} records the formal per-run limits.

\begin{table}[htbp]
\centering
\caption{Formal per-run solver budgets.}
\label{tab:budgets}
\begin{tabular}{lrrl}
\toprule
Suite & Time (s) & Memory (MB) & Stage-two allocation \\
\midrule
Validation & 600 & 2048 & up to 60 s \\
Nominal & 600 & 2048 & 540 s primary + up to 60 s \\
Two-stage stability & 600 & 2048 & profile-specific; see text \\
Factorial DOE & 600 & 2048 & 540 s primary + up to 60 s \\
SPBS & 600 & 2048 & 540 s primary + up to 60 s \\
Sensitivity & 600 & 2048 & 540 s primary + up to 60 s \\
Distribution shift & 1200 & 4096 & 1140 s primary + up to 60 s \\
\bottomrule
\end{tabular}
\end{table}

All three stability profiles have the same 600-s total and 2048-MB memory
budgets. Full and Quadratic-objective (linear-off) reserve up to 60 s for
linear and quadratic refinement, respectively; Stage2-off assigns all 600 s
to the primary solve. Their primary horizons are therefore unequal. Neither
PDI/time nor the reconstructed route metrics identify a pure stage-two causal
effect holding the stage-one incumbent and search horizon fixed. The matched
comparison instead measures each complete profile under a common total budget,
with route-wait and cadence outcomes reconstructed from saved event schedules
under identical definitions.

\section{Independent SPBS warm-start experiment}

Table~\ref{tab:spbs} reports the controlled SPBS comparison. Automatic SPBS
provides a valid incumbent in every run and reduces PDI by 7194.27 (26.19\%).
The matched instance-level effects on incumbent availability and PDI both
remain significant after multiplicity adjustment. Both conditions start their
600-s timed production solve after any SPBS construction has finished; hence
the comparison isolates the value of supplying the verified incumbent but
does not include its offline generation cost.

\begin{table}[htbp]
\centering
\caption{SPBS warm-start comparison over 48 instances and ten matched seeds.}
\label{tab:spbs}
\begin{tabular}{lrrrr}
\toprule
Condition & Runs & Incumbent (\%) & Optimal (\%) & Mean PDI \\
\midrule
Automatic SPBS & 480 & \textbf{100.00} & 54.58 & \textbf{20278.78} \\
No warm start & 480 & 51.67 & 51.67 & 27473.05 \\
\bottomrule
\end{tabular}
\end{table}

The incumbent-rate difference is 48.33 percentage points (95\% CI
$[34.58,62.29]$; Holm $p=2.55\times10^{-5}$). The PDI difference is 7194.27
(95\% CI $[5145.78,9317.61]$; Holm $p=3.46\times10^{-8}$). The optimal-rate
and elapsed-time differences do not remain significant after Holm adjustment.

\section{One-factor parameter sensitivity}

Table~\ref{tab:sensitivity} reports all eight prespecified settings. Each row
contains ten default-SCIP runs and every run provides an independently valid
incumbent. The experiment evaluates operational solvability under parameter
perturbations; because it contains one method, it is not used for a cross-
method robustness ranking.

\begin{table}[htbp]
\centering
\caption{One-factor sensitivity results for the 24-wafer $1{:}1$ mixed-flow case.}
\label{tab:sensitivity}
\begin{tabular}{lrrr}
\toprule
Setting & Optimal (\%) & Mean PDI & Mean time (s) \\
\midrule
Nominal & 20 & 40728.91 & 515.09 \\
Motion scale 0.8 & 40 & 37126.93 & 494.55 \\
Motion scale 1.2 & 0 & 42478.32 & 539.79 \\
Processing scale 0.8 & 60 & 33961.82 & 459.63 \\
Processing scale 1.2 & 10 & 38663.67 & 508.40 \\
Cleaning interval 3 & 10 & 38429.95 & 517.26 \\
Cleaning interval 10 & 40 & 28899.41 & 497.91 \\
Vacuum-zone dummy-wafer inventory 10 & 20 & 38152.52 & 484.80 \\
\bottomrule
\end{tabular}
\end{table}

\section{Distribution-shift stress test}

The OOD suite contains nine larger 40--64-wafer configurations. Its complete
729-run matrix has no duplicate key, execution error, solution-write error, or
invalid saved solution. Table~\ref{tab:ood} reports the instance-first
summaries. Every method attains 100\% incumbent availability and 0\% proven
optimality under the common horizon.

\begin{table}[htbp]
\centering
\caption{Complete OOD comparison over nine larger configurations.}
\label{tab:ood}
\begin{tabular}{lrrr}
\toprule
Method & Runs & Mean PDI & Mean gap (\%) \\
\midrule
SCIP & 27 & 73670.99 & 171.35 \\
Adaptive cut selection & 27 & \textbf{70915.54} & \textbf{160.78} \\
HEM-13D & 135 & 73003.59 & 165.88 \\
HEM-13D+Beam & 135 & 73165.58 & 167.04 \\
23D Features Only & 135 & 74187.22 & 169.96 \\
Structure+Greedy & 135 & 73507.76 & 168.54 \\
RDMCT-HBS & 135 & 73919.15 & 170.02 \\
\bottomrule
\end{tabular}
\end{table}

Adaptive cut selection has the lowest OOD mean PDI. RDMCT-HBS has a mean
terminal gap 0.78\% below SCIP. None of the OOD pairwise contrasts remains
significant after global Holm adjustment. The complete suite establishes the
scalable feasibility of the shared SPBS--MILP--SCIP pipeline and complements
the observed nominal-range mean anytime advantage reported in the main paper.

\section{Factorial-model diagnostics}

The factorial design contains 60 complete cells and five solver seeds per
cell. The prespecified PDI model has $R^2=0.9980$. Wafer count and recipe
composition exhibit strong interactions, whereas all prespecified processing-
scale main-effect and interaction confidence intervals include zero.
Figure~\ref{fig:diagnostics} provides the residual and influence diagnostics.

\begin{figure}[htbp]
\centering
\includegraphics[width=0.96\linewidth]{Figure_S1_DOE_diagnostics.png}
\caption{Diagnostic plots for the prespecified factorial PDI model fitted to
the complete 300-run design.}
\label{fig:diagnostics}
\end{figure}

\section{Reproducibility record}

The reproducibility package preserves the instance manifest, solver and
training seeds, method configuration, time and memory limits, trained-model
metadata, raw row records, saved incumbent solutions, independently derived
manufacturing metrics, and machine-readable statistical summaries. Every
formal suite is accepted only when its expected matrix is complete, duplicate
keys and execution errors are absent, saved solutions are covered one-to-one
by independent validation, and the statistical input count equals the frozen
raw count.

\end{document}
